% Source-derived chapter 07: Hyperk\"ahler aspects
% Original source: very-stable-higgs-bundles-equivariant-multiplicity-mirror-symmetry
\chapter{Hyperk\"ahler aspects}
\label{very-stable-higgs-bundles-equivariant-multiplicity-mirror-symmetry-chapter-07}
\source{very-stable-higgs-bundles-equivariant-multiplicity-mirror-symmetry}

\begin{remark}[Source section]
\label{very-stable-higgs-bundles-equivariant-multiplicity-mirror-symmetry-chapter-07-source}
\source{very-stable-higgs-bundles-equivariant-multiplicity-mirror-symmetry}
\source{very-stable-higgs-bundles-equivariant-multiplicity-mirror-symmetry}
This chapter reproduces the corresponding section of the retrieved
LaTeX source, including its definitions, statements, examples, and
proofs. It has not yet been translated into Lean.
\notready
\end{remark}

% The source section title was: Hyperk\"ahler aspects
\label{hyper}
\source{very-stable-higgs-bundles-equivariant-multiplicity-mirror-symmetry}



\subsection{Hyperholomorphic connections}

Thus far we have considered $\M$ as a variety with a holomorphic symplectic structure in which the upward flows are complex Lagrangian subvarieties. At this point we bring into play the hyperk\"ahler metric on $\M$. This is a metric which is K\"ahler with respect to complex structures $I,J,K$ which obey the quaternionic relations $I^2=J^2=K^2=IJK=-1$. It exists as a consequence of the theorem \cite{hitchin},\cite{simpson1},\cite{simpson1b} that a stable Higgs bundle admits a Hermitian metric, unique up to gauge equivalence, satisfying the equation $F_A+[\Phi,\Phi^*]=0$. 

Within this context, mirror symmetry predicts that a BAA-brane should correspond to a BBB-brane on the mirror of $\M$, which is to be understood as the statement that to a flat vector bundle on a complex Lagrangian on $\M$ there should exist a hyperholomorphic connection on a vector bundle over a hyperk\"ahler submanifold of the mirror.  A hyperholomorphic connection is a unitary connection whose curvature is of type $(1,1)$ with respect to each complex structure in the hyperk\"ahler family  of complex structures $x_1I+x_2J+x_3K$ with   $(x_1,x_2,x_3)\in S^2\subset \R^3$.  At this point in time, we have few constructions of hyperholomorphic connections on ${\mathcal M}$ so verifying the predictions is an important issue.

Since the fibration structure of $\M$ gives a natural setting for the Strominger-Yau-Zaslow approach to mirror symmetry, it is the Fourier-Mukai transform approach which can be applied. In the original SYZ context of a special Lagrangian fibration Arinkin and Polishchuk \cite{arinkin} showed how a real Lagrangian can describe a holomorphic bundle on the mirror. This could be applied to the current situation but would give at most  a hyperholomorphic connection relative to the {\emph semi-flat} hyperk\"ahler metric which is defined over the regular locus in $ \calA$ and is only an approximation to the actual metric. We have described the Fourier-Mukai transform of the trivial bundle over an upward flow Lagrangian in terms of universal bundles over the mirror moduli space, so to verify the prediction we need to construct a hyperholomorphic connection, using the genuine hyperk\"ahler metric, on a universal bundle  over the mirror of the moduli space. Since we are working with $GL_n(\C)$ Higgs bundles here, and this group is isomorphic to its Langlands dual, the mirror is the same moduli space. Twistor theory allows us to express this by holomorphic means using the direct analogy of the Atiyah-Ward description of instantons \cite{AHS}, \cite{Sal} -- a hyperholomorphic connection on $\R^4$ is precisely an  anti-self-dual connection. 

Recall that the twistor space $Z$ of a hyperk\"ahler manifold $M$ is  the manifold $M\times S^2$ endowed with a complex structure such that the fibre over $(x_1,x_2,x_3)\in S^2$ is $M$ with complex structure $x_1I+x_2J+x_3K$. The projection onto the second factor is a holomorphic fibration $p: Z\rightarrow \PP^1$ with a real structure $\sigma$ covering the antipodal map on $S^2$ and each 2-sphere $(m,S^2)$ is a real holomorphic section, called a twistor line.   Then a unitary hyperholomorphic connection on a vector bundle is equivalent to a holomorphic bundle on $Z$ which is trivial on each twistor line and real with respect to $\sigma$.

A description of the twistor space for Higgs bundles due to Deligne and Simpson \cite{simpson3} is phrased in terms of $\lambda$-connections. A holomorphic $\lambda$-connection on a vector bundle $E$ is a differential operator $D:E\rightarrow E\otimes K$ satisfying 
$D(fs)=fDs+\lambda s\otimes df$ for a local holomorphic section $s$ and function $f$. If $\lambda\ne 0$ then $\lambda^{-1}D$ is a holomorphic, and therefore flat, connection and for $\lambda=0$ this is a Higgs field. There is an algebraic construction of a moduli space ${\M}^s_{\mathrm{Hod}}$ of stable $\lambda$-connections and the parameter $\lambda$ gives a projection to $\C$. 

A solution to the Higgs bundle equations gives rise to a $\bar\partial$-operator $\nabla^{0,1}_A+\zeta\Phi^* =\bar\partial_A+\zeta\Phi^*$ parametrized by $\zeta\in \C$ which commutes with $\zeta\partial _A+ \Phi$, so defining a $\lambda$-connection with $\lambda=\zeta$ for the holomorphic structure on $E$ defined by  the $\bar\partial$-operator. This is a holomorphic section of ${\mathcal M}_{\mathrm{Hod}}\rightarrow \C$ and the twistor space ${\mathcal Z}$ is constructed so as to extend this to $\zeta=\infty\in \PP^1$ and make this a twistor line. It  is defined by  taking the moduli space of $\lambda$-connections  and the corresponding one for the conjugate complex structure on $C$, and identifying with the moduli space of  flat connections over $\zeta\in \C^*\subset \PP^1$ by means of the holonomy, the representation of $\pi_1(C)$ in $GL_n(\C)$, which is independent of any holomorphic structure. Then identifying $\lambda=\zeta$ on one side with $\lambda=\zeta^{-1}$ on the other gives ${\mathcal Z}$. 



When $\zeta\ne 0$ we have the flat connection
\begin{equation}
\nabla_A+\zeta^{-1}\Phi+\zeta \Phi^*.
\label{flatconn} 
\source{very-stable-higgs-bundles-equivariant-multiplicity-mirror-symmetry}
\end{equation}
and as $\zeta$  tends to $\infty$, $\partial_A+\zeta^{-1}\Phi$ defines a holomorphic structure with respect to the conjugate complex structure on $C$ and $\zeta^{-1}\bar\partial_A+\Phi^*$  a $\lambda$-connection for $\lambda=\zeta^{-1}$.   This defines a section of the fibration $p: {\mathcal Z}\rightarrow \PP^1$.

Since $A$ is a unitary connection the antilinear involution $a\mapsto -a^*$ on $\End(E)$ coupled with the antipodal map $\zeta\mapsto -1/\bar \zeta$ acting on $\PP^1$ transforms $\nabla^{1,0} _A+ \zeta^{-1}\Phi$ to  $\nabla^{0,1}_A+\bar \zeta \Phi^*$ and so a solution to the Higgs bundle equations  defines a real section, and this corresponds to a point of $\M$. 



\subsection{The connection on a universal bundle}
As described above, the Fourier-Mukai transform of the upward flow of a very stable fixed point  of type $(1,1,\dots,1)$ can be expressed as a tensor product of exterior powers of the universal bundle corresponding to  points $c\in C$. A hyperholomorphic connection on the universal bundle will then induce one on the associated tensor or exterior powers. 

In the  complex structure  $\zeta \ne 0,\infty$ there is a straightforward description of  the universal bundle: given a framing of a bundle $E$ at $c$, the holonomy of a flat connection on $E$ around closed curves with base point $c$ defines a homomorphism from $\pi_1(C,c)$ to $\GL_n(\C)$.   We denote $\calR_\B=\Hom(\pi_1(C,c)\to \GL_n(\C))$, and the open locus of irreducible representations $\calR^s_\B$. We have the action of $\GL_n(\C)$ on $\calR_\B$ by conjugation, and denote the affine quotient by $\M_\B$, the character variety. Over the open subset $\M^s_\B=\calR^s_\B/\GL_n(\C)\subset \M_\B$ the  variety $\calR^s_\B$ is a principal $\PGL_n(\C)$-bundle, representing irreducible local systems framed at the point $c\in C$.

As described in Section 6, there is an obstruction in  $H^2({\M^s_\B}, \Z_n)$ to lifting to a  principal $\GL_n(\C)$-bundle. Given local liftings over  open sets $U_i$, the obstruction class is represented by a flat line bundle $L_{ij}$ over $U_i\cap U_j$ with $L_{ij}^n$ trivial -- the data for a gerbe $\theta$. If the universal vector bundle ${\mathbb E}$  does not exist globally, it  does over the neighbourhoods $U_i$, and a connection is a locally defined notion. Moreover tensoring with a flat line bundle $L_{ij}$ acts trivially on the curvature and so takes a local hyperholomorphic connection to another one, so the notion of hyperholomorphic passes over to the situation where we need to twist by a flat gerbe.

For the complex structure $\zeta=0$ or $\infty$, the moduli space of framed Higgs bundles was constructed by Simpson in    \cite[Theorem 4.10]{simpson1}, and in \cite[Proposition 4.1]{simpson3} he  outlined, by the same approach, a construction for $\lambda$-connections. We denote by $\calR^s_\Hod$ the framed moduli space of stable $\lambda$-connections. This carries a $\GL_n(\C)$-action on the framing, so that it becomes a  principal $\PGL_n(\C)$-bundle over $\M^s_\Hod$.  Over $\lambda=0$ this principal bundle $\calR^s_\Hod\to \M^s_\Hod$ restricts to $\calR^s_\Dol\to \M_\Dol^s=\M^s$  the framed moduli space of stable rank $n$ degree $0$ Higgs bundles, which is a $\PGL_n(\C)$-principal bundle over $\M^s$. 

When $\lambda\ne 0$ we have the complex analytically equivalent alternative description $\calR^s_\B\to \M^s_\B$ of the framed moduli space in terms of the character variety, using the holonomy of the flat connection. This  is clearly compatible with the identification in the construction of  the twistor space ${\mathcal Z}$,  so we have a holomorphic $\PGL_n(\C)$-bundle over ${\mathcal Z}$. The final point is to show this is trivial on a real twistor line. 

A real twistor line is defined by  a  solution to the Higgs bundle equations which, as $\zeta$ varies,  is a pair $\bar\partial_A+\zeta\Phi^*$,  $\zeta\partial _A+ \Phi$ on the {\emph {same}} $C^{\infty}$ vector bundle $E$. A basis $e_1,\dots, e_n$ of the fibre $E_c$ then defines a framing for all holomorphic structures and we need this to be holomorphic in $\zeta$. But  from the integrability theorem of the holomorphic structure $\bar\partial_A+\zeta\Phi^*$ we can find a basis of local sections $s_1,\dots, s_n$ in a neighbourhood of $c$   with $(\bar\partial_A+\zeta\Phi^*)s_i=0$ and  these vary holomorphically with $\zeta$. 
Then evaluation at $c$ gives a frame which is holomorphic in $\zeta$ and similarly using the conjugate complex structure for $\zeta\ne 0$. . 




\begin{remark}
\source{very-stable-higgs-bundles-equivariant-multiplicity-mirror-symmetry}
\notready 
	
	A formal explanation for the  existence of the hyperholomorphic connection is provided by the infinite-dimensional hyperk\"ahler quotient construction of ${\mathcal M}$ itself \cite{hitchin}. This considers the action of the group ${\mathcal G}$ of unitary gauge transformations on the infinite-dimensional flat hyperk\"ahler manifold which is the cotangent bundle of the space ${\mathfrak A}$ of $\bar\partial$-operators $\bar\partial_A$ on the underlying $C^{\infty}$ vector bundle $E$. The zero set of the hyperk\"ahler moment map for this action consists of pairs $(A,\Phi)$ satisfying the Higgs bundle equations. This space is formally a principal ${\mathcal G}/U(1)$ bundle over the quotient ${\mathcal M}$ and the restriction of the flat $L^2$ metric on $T^*{\mathfrak A}$ defines a distribution orthogonal to the orbit directions, and this is a connection on the principal bundle. In any hyperk\"ahler quotient this is hyperholomorphic. Now given a framing at $c$, the evaluation map gives a homomorphism from ${\mathcal G}/U(1)$ to $PU(n)$ and an associated hyperholomorphic connection on the universal bundle.
\end{remark}






%A universal bundle for complex structure $I$, the moduli space of Higgs bundles, was constructed on ${\mathcal M}\times C$ in \cite{hausel-vanishing}. Fixing a point $x\in  C$ and a framing at $x$ gives the principal bundle under consideration here. 




Thus we have proved the following.

\begin{theorem}
\source{very-stable-higgs-bundles-equivariant-multiplicity-mirror-symmetry}
\notready There is a hyperholomorphic connection on the $\PGL_n(\C)$ bundle $\calR^s_\Dol$ on $\M^s$, which is associated to the projective bundle $\P(\E_c)$. Moreover, there is a hyperholomorphic structure on the  vector bundle $\E_c$ on $\M^s$ twisted by the gerbe  $\theta$. 	
\end{theorem}

\subsection{The universal Higgs field} 
For each point $c\in C$ we obtain the universal bundle $\E_c$ with a hyperholomorphic connection on ${\mathcal M}$, and so $C$ itself is the parameter space for  a family of such bundles with connection.  Equivalently, these correspond to holomorphic bundles on the twistor space ${\mathcal Z}$. In particular, considering first order deformations, there is a natural map from the tangent space of $C$ at $c$ to $H^1({\mathcal Z},\End (\E_c))$. This space can be interpreted  in terms of variations of hyperholomorphic connections  by using the Penrose transform in the current quaternionic context \cite{Sal}. It is isomorphic to the first cohomology group of an elliptic complex generalizing  the deformation complex for anti-self-dual connections in four dimensions in \cite{AHS}. 

Since $\M$ is the moduli space of pairs $(E,\Phi)$ the universal bundle should have a universal Higgs field to accompany it and indeed a global construction of the universal pair  can be found in  \cite{hausel-vanishing} for the $(d,n)=1$ case
and  \cite[Theorem 4.5]{simpson1}  for a local construction in general. 
Given the role of the universal bundle, the reader may wonder how the universal Higgs field at $c$ enters the picture. We shall see here that it defines the infinitesimal variation of the hyperholomorphic connection on the universal bundle.


The starting point is the observation that the space $\Hom(\pi_1(C,c), \PGL_n(\C))$ is independent of $c$. More accurately, choosing a path from $c$ to $y$, there is a biholomorphism between the two spaces which only depends on the homotopy class of the path. In our situation it is the holonomy of the flat connection along that path which identifies the two. We have a flat connection when $\zeta \ne 0$ or $\infty$ and this suggests that the deformation class in $H^1({\mathcal Z},\End (\E_c))$ is supported on the fibres $Z_0,Z_{\infty}\subset {\mathcal Z}$ over these two points in $\PP^1$.

Given a tangent vector  $X\in T_c(C)$, we can contract with the  Higgs field $\Phi_c\in (\End (E) \otimes K)_c$ to obtain $\Phi(X)$ which is a section of the universal bundle $\End (\E_c)$ over  the moduli space of Higgs bundles  ${\mathcal M}$ which we now think of as the fibre $Z_0$ over $\zeta=0$ in the twistor space. Similarly $-\Phi^*$ defines a holomorphic section $-\Phi^*(X)$ over $Z_{\infty}$, the moduli space  for the conjugate complex structure. Let $s$ be the pullback  to ${\mathcal Z}\rightarrow \PP^1$ of the holomorphic section $\zeta d/d\zeta$ of $T\PP^1\cong {\mathcal O}(2)$ on $\PP^1$, then we have an exact sequence of sheaves on  ${\mathcal Z}$:
$$0\rightarrow {\mathcal O}(\End (\E_c))\stackrel {s} \rightarrow {\mathcal O}(\End (\E_c)(2))\rightarrow {\mathcal O}_{Z_0+Z_{\infty}}(\End (\E_c)(2))\rightarrow 0.$$
The vector field  $d/d\zeta$ trivializes ${\mathcal O}(2)$ at $\zeta=0$ and similarly at infinity so the pair $(\Phi(X),-\Phi^*(X))$ may be regarded as a section of $\End (\E_c)(2)$ over $Z_0+Z_{\infty}$. Then we have: 
\begin{proposition}
\source{very-stable-higgs-bundles-equivariant-multiplicity-mirror-symmetry}
\notready The map $TC_c\rightarrow H^1({\mathcal Z},\End (\E_c))$ defining the first order deformation of the universal bundle at $c$ is $\delta(\Phi(X),-\Phi^*(X))$ where $\Phi$ is the universal Higgs field at $c$ and  $\delta: H^0({Z_0+Z_{\infty}},\End (\E_c)(2))\rightarrow H^1({\mathcal Z},\End (\E_c))$   is the connecting homomorphism in the cohomology of the above exact sequence. 
\end{proposition} 
\begin{proof} For $y\in C$ in a neighbourhood of $c$,  parallel translation of a flat connection  takes the principal bundle $\Hom(\pi_1(C,c), \PGL_n(\C))$ to $\Hom(\pi_1(C,y), \PGL_n(\C))$. This is an isomorphism  of the universal bundle at $c$ to the universal bundle at $y$, and to first order at $c$ it is the horizontal lift of a tangent vector $X$ to the universal bundle over $\M\times C$, the lift defined by the connection. This is a trivial deformation of holomorphic  principal bundles where $\M$ has the complex structure of the moduli space of flat connections, so since $\zeta\ne 0,\infty$ describes the points of ${\mathcal Z}$ via flat connections, it gives a trivialization of the deformation of the universal bundle on this part of ${\mathcal Z}$. The flat connection is $\partial _A+ \Phi/\zeta$ and so the trivialization does not extend to $\zeta=0$. 
	
	We may consider the first order deformation of the universal bundle as the Kodaira-Spencer class for the deformation of complex structure of the principal bundle ${\mathcal P}_c\rightarrow {\mathcal Z}$ as $c\in C$ varies. 
	In \v Cech cohomology, the deformation class in $H^1({\mathcal Z},\End (\E_c))\subset H^1({\mathcal P}_c, T)$ is given using an open covering  $\{U_i\}$, taking  local lifts $X_i$ of the tangent vector $X$ and defining the cocycle $X_j-X_i$, a section of $\End (\E_c)$ over $U_i\cap U_j$. In our case, for $\zeta\ne 0$, there is a global lift given by the horizontal space of the connection $\partial _A+ \Phi/\zeta$. If we multiply by $\zeta$ then this is global for all $\zeta$. Repeating this process near $\zeta=\infty$ shows that the class maps to zero in $H^1({\mathcal Z},(\End (\E_c))\stackrel {s} \rightarrow H^1({\mathcal Z},\End (\E_c)(2))$ and consequently is in the image of the coboundary map $\delta:
	H^0({Z_0+Z_{\infty}},(\End (\E_c)(2))\rightarrow H^1({\mathcal Z},(\End (\E_c))$.
	
	But the polar part of $\partial _A+ \Phi/\zeta$ at $\zeta=0$ is $\Phi(X)$ and 
	we see that the class is given by extending $\Phi(X)$ and $-\Phi^*(X)$ as sections of $\End (\E_c)(2)$ and dividing by $s$, which is the coboundary map  in \v{C}ech cohomology.
\end{proof}
